\chapter{The Boussinesq Approximation}

Sounds waves are usually among the fastest frequencies in the flow. In
Earth's atmosphere, a sound wave propagates at about 300 m/s, or over
1000 km/h. Say you want to model the atmosphere in a box over
1000 kilometers across, with a typical resolution of 10\,km. Your
CFL condition, with Courant number 0.5, yields a timestep
of $\Delta t = C\Delta x / u  = 15 s$. While you can get away with it
for simulating weather over a week, this high frequency would grind a
simulation to a halt if you were interested in modeling, say, climate,
over years or decades. 

We could eliminate the sound waves by assuming incompressibility. Yet
this would also remove any density variation, and thus buoyancy-driven
flows such as convection. From experience at Earth's atmosphere, or
simply boiling water in a kettle, we know that near incompressible
fluids allow for convection. So full incompressibily is not a good
approximation if we are interested in buyoancy effects. 

Enter the Boussinesq approximation. This approximation removes sound
waves, while allowing for buoyancy-driven flows by retaining density
variations only when the term is multiplied by gravity.

\subsection{Decomposition}

We will use the following notation for
base state (zero subscript), stratification (hat), perturbation
(prime), base+stratification (tilde), and stratification plus
perturbation (delta)

\begin{eqnarray}
  a(\v{x},t) &=& a_0 + \delta a(\v{x},t)\\
&=&a_0 + \hat{a}(z) + a^\prime(\v{x},t)\\
&=&\tilde{a}(z) + a^\prime(\v{x},t)
\end{eqnarray}

%\noindent which means
%
%\begin{eqnarray}
%\delta a(\v{x},t)&=&\hat{a}(z) + a^\prime(\v{x},t)\\
%\tilde{a}(z) &=& a_0 + \hat{a}(z)
%\end{eqnarray}

\subsection{Momentum Equation}

Consider the compressible momentum equation

\beq
\frac{D u_z}{Dt}=-\frac{1}{\rho}\frac{\partial p}{\partial z} -g
\eeq

\noindent where $g>0$.  Multiply by density, splitting also 
density and pressure into base and stratification $\tilde{p}(z)$ plus
pertubation $p^\prime(\v{x},t)$. 

\begin{equation}
(\rho_0+\delta\rho) \frac{Du_z}{Dt}=-\frac{\partial
  \tilde{p}(z)}{\partial z} -\frac{\partial p^\prime}{\partial z} - g\rho_0- g\delta\rho
\end{equation}

The base state is hydrostatic, canceling $\partial_z\tilde{p}$ and
$g\rho_0$. Considering $\delta\rho \ll \rho_0$, we have 

\begin{equation}
\rho_0\frac{Du_z}{Dt}= -\frac{\partial p^\prime}{\partial z} - g\delta\rho
\end{equation}

\noindent defining enthalpy $d\Pi\equiv \rho_0^{-1}dp$ and also definining 
\beq
\Theta \equiv \frac{\delta\rho}{\rho_0}
\eeq

The momentum equation becomes 

\begin{equation}
\frac{Du_z}{Dt}= -\frac{\partial\Pi}{\partial z} -g\Theta
\end{equation}

Notice that $\Theta$ is dimensionless. We call it simply ``the buoyancy quantity''.

\subsection{Buoyancy Equation}

The buoyancy equation comes from the left hand side of the continuity equation 

\begin{equation}
\frac{D\delta\rho}{Dt} = 0 
\end{equation}

\noindent and expanding delta into tilde (base + stratification) and prime (perturbation)

\begin{eqnarray}
0 &=&\frac{D\left[\tilde\rho(z) + \rho^\prime(\v{x},t)\right]}{Dt} \\
      &=&u_z \partial_z \tilde\rho(z)  +\frac{D\rho^\prime(\v{x},t)}{Dt}
\end{eqnarray}

\noindent Multiplying by $1/\rho_0$ we have the equation for the buyoancy variable  

\beq
\frac{D\Theta^\prime}{Dt} = -u_z \partial_z \tilde{\Theta}
\eeq

\noindent We have the three equations, plus the trivial equation for the
velocity $u_x$ (and/or $u_y$) in the non-stratified direction(s) 

\beqn
\frac{Du_z}{Dt}&=& -\frac{\partial\Pi}{\partial z} -g\Theta\\
 \frac{Du_x}{Dt} &=& -\frac{\partial\Pi}{\partial x}\\
 \frac{D\Theta^\prime}{Dt} &=& -u_z \partial_z \tilde{\Theta}\\
\Div{\v{u}}&=&0
 \eeqn

\subsubsection{Brunt-V\"ais\"al\"a frequency and gravity waves}
 
Ignoring enthalphy, we can see that this system supports oscillations
in the vertical direction, which we recognize by solving the perturbation equations

 
\beqn
\frac{\partial^2  u_z^\prime}{\partial t^2}= g \partial_z \tilde{\Theta} u_z^\prime \\
\eeqn

\noindent which has solution $e^{-iN t}$ with $N = \sqrt{-g\partial_z
  \tilde{\Theta}}$. These waves are buoyancy waves, also called
internal gravity waves, as the restoring force is gravity. The
quantity $N$ is a frequency, called Brunt-V\"ais\"al\"a frequency.


So the equations are

\begin{empheq}[box=\fbox]{align}
\frac{Du_z}{Dt}&= -\frac{\partial\Pi}{\partial z} - g\Theta\\
 \frac{Du_x}{Dt} &= -\frac{\partial\Pi}{\partial x}\\
 \frac{D\Theta^\prime}{Dt} &= \frac{1}{g} N^2 u_z \\
\Div{\v{u}}&=0
\end{empheq}

Notice that the term $\partial_z \tilde{\Theta}$ is the vertical
stratification, which we write as $1/\Lambda$, so

\begin{equation}
  g=\Lambda N^2
\end{equation}

\noindent  The buoyancy and momentum equations are

\beqn
\frac{D\Theta}{Dt} &=& \frac{1}{\Lambda} u_z \\
\frac{Du_z}{Dt}&=& -\frac{\partial\Pi}{\partial z} -\Lambda N^2\Theta
\eeqn

\noindent we can absorb the stratification length $\Lambda$ into $\Theta$,
writing $\theta = \Lambda\Theta$, so finally 

\begin{empheq}[box=\fbox]{align}
\frac{D\theta}{Dt} &=  u_z \\
\frac{Du_z}{Dt}&= -\frac{\partial\Pi}{\partial z} - N^2\theta
\end{empheq}

$\theta$ is called potential temperature, as the density variations
are associated with temperature variations. 


\section{Shallow water}

Depth-integrated Navier-Stokes. 

Applies when $L\gg z$, $v_z \ll v$, and $dp/dz$ is hydrostatic.

The horizontal fluid velocity is constant through the fluid depth.

$Lz/Lz \ll 1$, dynamic motions are hydrostatic. 


timescale for vertical oscillation: $T \approx \frac{L_x}{c} \approx
\frac{L_x}{\sqrt{gL_z}}$

free-fall time: $T_{\rm ff}=\sqrt{L_z/g}$


$T_{\rm ff}/T = \sqrt{\frac{L_z}{g}} \sqrt{\frac{gL_z}{L_z}} =
\frac{L_z}{L_x} \ll 1 $


oscillation is much longer than free fall. The system can always
equilibrate, the dynamics is hydrostatic. 

\beq
\frac{\partial p}{\partial z} = -\rho g 
\eeq

\beq
dp = -\rho g dz
\eeq


\beq
p = -\int_{-h}^0 \rho g dz = h(x)
\eeq


for $z = L_z$, $p=-\rho g z = \rho g L_z$. 

$p/\rho$ has dimension of velocity square, so we can write $p=\rho c^2$,
with $c=\sqrt{gL_z}$.

Write the NS equations for incompressible 2D x-z flow 

\beqn
\frac{\partial u_x}{\partial t} + u_x \frac{\partial u_x}{\partial x}+ u_z\frac{\partial u_x}{\partial z} &=& -\frac{1}{\rho}\frac{\partial P}{\partial x}\\
\frac{\partial u_z}{\partial t} + u_x \frac{\partial u_z}{\partial x}+u_z\frac{\partial u_z}{\partial z} &=& -\frac{1}{\rho}\frac{\partial P}{\partial z} -g \\
\frac{\partial u_x}{\partial x} + \frac{\partial u_z}{\partial z} &=&0
\eeqn

From the last equation, the continuity equation

\beq
\frac{u_x}{L_x} \sim \frac{u_z}{L_z} \longrightarrow \frac{u_z}{u_x}
\sim \frac{L_z}{L_x} \ll 1 
\eeq

so we conclude that vertical velocities are negligible $u_z \ll u_x$,
or $u_z / u_x = \mathcal{O}(\epsilon)$, so $u_z \sim \epsilon u_x$.

Now let us consider the vertical velocity equation

\beq
\frac{\partial u_z}{\partial t} + u_x \frac{\partial u_z}{\partial x}+u_z\frac{\partial u_z}{\partial z} = -\frac{1}{\rho}\frac{\partial P}{\partial z} -g 
\eeq

and analyse the magnitude of the terms. The timescale is $T=L_x/c_g$,
so the terms are

\beqn
\frac{u_z}{T} &\sim& \frac{\epsilon u_x}{T} \sim \frac{\epsilon  c_g}{T} \sim \frac{\epsilon c_g^2}{L_x}\\
u_x \frac{\partial u_z}{\partial x}&\sim&\frac{u_xu_z}{L_x}\sim\frac{\epsilon c_g^2}{L_x}\\
u_z \frac{\partial u_z}{\partial z}&\sim&\frac{u_z^2}{L_z}\sim\frac{\epsilon c_g^2}{L_x}\\
-\frac{1}{\rho}\frac{\partial p}{\partial z} &\sim&\frac{c_g^2}{L_z}\sim\frac{c_g^2}{\epsilon L_x}
\eeqn

so the pressure term dominates, and indeed the velocities remain near
zero, and hydrostatic.The vertical NS equation reduces to

\beq
\frac{\partial p}{\partial z} = -\rho g
\label{eq:hydroz}
\eeq

Integrating the hydrostatic equilibrium equation

\beq
p = -\int_\eta^z \rho g dz^\prime = \rho g \left(\eta-z\right)
\label{eq:pressure-shallow}
\eeq

so

\beq
p(x,z,t) = \rho g \left[ \eta(x,t)-z\right]
\eeq

The fluid adjusts to restore vertical hydrostatic equilibrium as $\eta$
varies.

Now use this equation for P in the horizontal momentum equation

\beq
\frac{\partial u_x}{\partial t} + u_x \frac{\partial u_x}{\partial x}+u_z\frac{\partial u_x}{\partial z} = -\frac{1}{\rho}\frac{\partial p}{\partial x} 
\eeq

Substituting \eq{eq:pressure-shallow}, 

\beq
\frac{\partial p}{\partial x} = \rho g \frac{\partial \eta}{\partial x}
\eeq

and now using that to substitute the pressure gradient

\beq
\frac{\partial u_x}{\partial t} + u_x \frac{\partial u_x}{\partial x}+u_z\frac{\partial u_x}{\partial z} = -g \frac{\partial \eta}{\partial x}
\eeq

Notice that neither $g$ nor $\partial \eta$ depend on z. If we also
have that $\partial_z
u_x$ is zero, we find 

\beq
\frac{\partial u_x}{\partial t} + u_x \frac{\partial u_x}{\partial x} = -g \frac{\partial \eta}{\partial x}
\label{eq:ux-shallow}
\eeq

In Lagrangian form we can write it as

\beq
\frac{Du_x}{Dt} = -g \frac{\partial \eta}{\partial x}
\eeq

which is independent of $z$. So, if $\partial_z u_x=0$, then
$\frac{Du_x}{Dz}=0$. \classcomm{Not sure what is the significance of
  this argument}.


Consider now the continuity equation

\beq
\frac{\partial u_z}{\partial z} = -\frac{\partial u_x}{\partial x} 
\eeq

integrating

\beq
u_z = -z\frac{\partial u_x}{\partial x}+C(x,t) 
\label{eq:uz-shallow}
\eeq

now let us find the constant. With the boundary conditions, $z=-L_b$;
$u_z = - u_x\partial_x L_b$ \classcomm{where?}, so 

\beq
- u_x\frac{\partial L_b}{\partial x} = L_b\frac{\partial u}{\partial x}+C
\eeq

so

\beq
C = - u_x\frac{\partial L_b}{\partial x} - L_b\frac{\partial u}{\partial x}
\label{eq:const-shallow}
\eeq

substituting \eq{eq:const-shallow} into \eq{eq:uz-shallow}, we find

\beq
u_z = -\left(z+L_b\right)\frac{\partial u_x}{\partial x}-u_x\frac{\partial L_b}{\partial x}
\label{eq:uz2-shallow}
\eeq


consider now that z is the height of an upper moving boundary,
$z=\eta(x,t)$. Taking the time derivative on both sides,

\beq
u_z =  \frac{\partial \eta}{\partial t} + u_x\frac{\partial \eta}{\partial x}
\label{eq:advect-uz}
\eeq

equating \label{eq:advect-uz} with \label{eq:uz2-shallow}

\beq
 \frac{\partial \eta}{\partial t} + u_x\frac{\partial \eta}{\partial x} = -\left(z+L_b\right)\frac{\partial u_x}{\partial x}-u_x\frac{\partial L_b}{\partial x}
\eeq

define now the height $h\equiv \eta + L_b$

\beq
 \frac{\partial \eta}{\partial t} + u_x\frac{\partial }{\partial x} \left(\eta+L_b\right)= -\left(z+L_b\right)\frac{\partial u_x}{\partial x}
\eeq

since $L_b$ is time-independent, we add it to the first term and find

\beq
 \frac{\partial h}{\partial t} + u_x\frac{\partial h}{\partial x} = -h\frac{\partial u_x}{\partial x}
\label{eq:shallow-continuity}
\eeq

Now substituting $\eta=h-L_b$ on \eq{eq:ux-shallow}


\beq
\frac{\partial u_x}{\partial t} + u_x \frac{\partial u_x}{\partial x}
= - g \frac{\partial h}{\partial x} + g \frac{\partial L_b}{\partial x}
\label{eq:ux-shallow-final}
\eeq

the last term in the RHS is a forcing -- a slope. The first term in
the RHS behaves like a pressure. 

\eq{eq:shallow-continuity} and \eq{eq:ux-shallow-final} are the
shallow water equations. Notice that indeed they allow for
oscillations:

\beq
\frac{\partial^2 h^\prime}{\partial t^2} = gh \frac{\partial^2 h^\prime}{\partial x^2}
\eeq

so the surface of the fluid has perturbations that propagate with
velocity $c_g=\sqrt{gh}$. These are {\it surface gravity waves}.

The analysis considered one non-z dimension (x), but there is symmetry
with respect to x and y, so the same result can be easily generalized
to 2D, yielding


\beqn
 \frac{\partial h}{\partial t} + \left(\v{u}\cdot \del\right) h &=& -h\Div{\v{u}}\\
\frac{\partial \v{u}}{\partial t} + \left(\v{u}\cdot \del\right) \v{u}&=& - g \grad{h} + g \grad{L_b}
\eeqn

Where the gradients and vectors are supposed xy (2D) only. 

\subsection{Dispersion of surface gravity waves - solitons}

Notice how gravity waves are {\it non-dispersive}. All the
wavenumbers travel at the same speed (phase velocity $\omega/k$ and
group velocity $d\omega/dk$ are equal), meaning the wave's shape
remains constant as it propagates. Surface gravity waves can form
solitons (solitary waves) if the weak nonlinearity is perfectly
balanced by the weak dispersion.

Notice that the dispersion for deep water waves is dispersive. The
dispersion relation is $\omega = \sqrt{gk}$, so the phase velocity is
$c=\omega/k = \sqrt{g/k}$, whereas the group velocity is $d\omega/dk =
\frac{1}{2}\sqrt{g/k}$, or half the phase velocity.  The phase
velocity being $\sqrt{g/k}$ means that longer wavelength waves travel
faster than shorter ones in deep water. 

\subsection{Beach slope solution}

The constant beach slope solution follows from the shallow water equations by assuming a linearly varying depth 

\beq
h=Sx
\eeq

This leads to a steady solution where velocity increases as depth decreases, and the free surface follows a modified version of the depth profile.

\subsection{The atmosphere as shallow water}

\classcomm{Show Ali's work}


\section{Convection}

%\subsection{Convection in high Reynolds number. Schwarzschild criterion}

\subsection{Rayleigh number and Prandl number}

\subsection{Convection in low Reynolds number. Mantle convection}

Suppose that a fluid is held between two large, flat, parallel plates,
as shown in Figure).The lower plate is maintained at temperature T= T0
+ T and the upper plate
at temperature T0. At low values of T the fluid remains stagnant and
heat passes from one plate to the other by molecular conduction. Of course, there is an upward buoyancy force which is
greatest near the lower plate and tends to drive the fluid upward. However, at low values of T
this force is exactly balanced by a vertical pressure gradient. We now slowly increase T. At a
critical value of T the fluid suddenly starts to convect as shown in
Figure. 

The flow consists of hot rising fluid and cold falling regions. This takes the form of regular
cells, called Bénard cells, which are reminiscent of Taylor vortices.2 The rising fluid near the top
of the layer loses its heat to the upper plate by thermal conduction, whereupon it starts to fall.
As this cold fluid approaches the lower plate it starts to heat up and sooner or later it reverses
direction and starts to rise again. Thus we have a cycle in which
potential energy is continuously
released as light fluid rises and dense fluid falls. Under steady conditions the rate of working of
the buoyancy force is exactly balanced by viscous dissipation within
the fluid.
The transition from the quiescent state to an array of steady convection cells is, of course,
triggered by an instability associated with the buoyancy force. If we perturb the quiescent state,
allowing hot fluid to rise and cold fluid to fall, then potential energy will be released and, if the
viscous forces are not too excessive, the fluid will accelerate, giving rise to an instability. On the
other hand, if the viscous forces are large then viscous dissipation
can destroy all of the potential
energy released by the perturbation. In this case the quiescent configuration is stable. Thus we
expect the quiescent state to be stable if T is small and $\nu$ large, and
unstable if T is large and $\nu$ is small.

Now suppose Ra is slowly increased above its critical value. Then
there is a point at which the Bénard cells themselves become unstable,
and more complex, unsteady structures start to appear.
Eventually, when Ra is large enough, we reach a state of turbulent convection.
 
It seems that the general picture that emerges from Taylor’s experiments is
the following. For high viscosities the basic configuration is stable. As the viscosity is reduced we
soon reach a point where the basic equilibrium is unstable to infinitesimal perturbations and the
system bifurcates (changes) to a more complex state, consisting of a steady laminar flow in the
form of regular cells. As the viscosity is reduced even further the new flow itself becomes unstable
and we arrive at a more complex motion. Subsequent reductions in viscosity give rise to
progressively more complex flows until eventually,for sufficiently small $\nu$, the flow is fully turbulent.
At this point the flow field consists of a mean(time-averaged) component plus a random, chaotic
motion.This sequence of stes is called the {\it transition to turbulence}.


